\documentclass[11pt, a4paper]{article}

% ── Encoding & Language ───────────────────────────────────────────────────────
\usepackage[utf8]{inputenc}
\usepackage[T1]{fontenc}
\usepackage[ngerman]{babel}

% ── Layout ───────────────────────────────────────────────────────────────────
\usepackage[top=2.5cm, bottom=2.5cm, left=3cm, right=2.5cm]{geometry}
\usepackage{setspace}
\onehalfspacing
\setlength{\parindent}{0pt}
\setlength{\parskip}{5pt}
\usepackage{lmodern}

% ── Math ─────────────────────────────────────────────────────────────────────
\usepackage{amsmath, amssymb}

% ── Tables & Figures ─────────────────────────────────────────────────────────
\usepackage{booktabs}
\usepackage{tabularx}
\usepackage{float}
\usepackage{graphicx}
\usepackage[labelfont=bf, skip=4pt]{caption}
\graphicspath{{figures/}{figures/02_sentiment_analysis/}{figures/03_correlation_modeling/}}

% ── Colors & Links ───────────────────────────────────────────────────────────
\usepackage[dvipsnames]{xcolor}
\definecolor{linkblue}{HTML}{1A5276}
\definecolor{codebg}{HTML}{F5F5F5}
\usepackage[
  colorlinks=true,
  linkcolor=linkblue,
  citecolor=linkblue,
  urlcolor=linkblue,
  pdftitle={Technischer Methodenbericht: WebMiningProject},
]{hyperref}

% ── Code listings ────────────────────────────────────────────────────────────
\usepackage{listings}
\lstset{
  basicstyle=\ttfamily\small,
  breaklines=true,
  frame=single,
  backgroundcolor=\color{codebg},
  keywordstyle=\color{blue!70!black},
  commentstyle=\color{gray},
  stringstyle=\color{orange!80!black},
  showstringspaces=false,
  numbers=left,
  numberstyle=\tiny\color{gray},
  numbersep=6pt,
  literate=
    {ä}{{\"{a}}}1 {ö}{{\"{o}}}1 {ü}{{\"{u}}}1
    {Ä}{{\"{A}}}1 {Ö}{{\"{O}}}1 {Ü}{{\"{U}}}1 {ß}{{\ss}}1,
}

% ── Bibliography ─────────────────────────────────────────────────────────────
\usepackage[
  backend=biber,
  style=authoryear,
  sorting=nyt,
  maxnames=3,
]{biblatex}
\addbibresource{references.bib}

% ── Misc ─────────────────────────────────────────────────────────────────────
\usepackage{enumitem}
\usepackage{microtype}
\usepackage{csquotes}
\usepackage{titlesec}
\titlespacing*{\section}{0pt}{14pt}{6pt}
\titlespacing*{\subsection}{0pt}{10pt}{4pt}

% =============================================================================
\begin{document}
% =============================================================================

\begin{titlepage}
  \centering
  \vspace*{2cm}
  {\LARGE\bfseries Technischer Methodenbericht\par}
  \vspace{0.8cm}
  {\large Web Mining: Sentiment-Analyse von X\,(Twitter) und Aktienkursrenditen\par}
  \vspace{0.6cm}
  {\normalsize Softwarearchitektur und Implementierungsdetails\par}
  \vspace{2cm}
  {\large Ilker\par}
  \vspace{0.4cm}
  {\normalsize M.\,Sc. Data Science\par}
  \vfill
  {\normalsize \today\par}
\end{titlepage}

\tableofcontents
\newpage

% =============================================================================
\section{Einleitung}
\label{sec:einleitung}
% =============================================================================

Dieser Bericht dokumentiert die softwaretechnische Umsetzung des WebMiningProject.
Im Vordergrund stehen Architekturentscheidungen, Implementierungsdetails und das
Zusammenspiel der Systemkomponenten. Der statistische Analyseteil wird nur insoweit
beschrieben, als er konkrete Implementierungsentscheidungen begründet; inhaltliche
Interpretation der Ergebnisse findet sich im Hauptbericht.

Das System gliedert sich in vier Hauptbereiche:
(1)~Datenerhebung via X\,(Twitter)-Crawler,
(2)~Kursdaten-Pipeline über Yahoo Finance,
(3)~Sentiment-Analyse mit einem vortrainierten Transformer-Modell sowie
(4)~statistische Korrelationsmodellierung in Jupyter-Notebooks.
Die gesamte Implementierung erfolgte in Python~3.12.3 mit einer virtuellen
Umgebung (\texttt{.venv}).

% =============================================================================
\section{Systemarchitektur}
\label{sec:architektur}
% =============================================================================

\subsection{Verzeichnisstruktur}

Das Projekt folgt einer strikten Trennung zwischen Quellcode (\texttt{src/}),
Analysenotebooks (\texttt{notebooks/}), Rohdaten (\texttt{data/raw/}) und
verarbeiteten Daten (\texttt{data/processed/}):

\begin{lstlisting}[language={}, caption={Projektstruktur}]
WebMiningProject/
|-- src/
|   |-- scraping/       # X/Twitter-Crawler (Playwright)
|   |-- finance/        # Kursdaten-Pipeline (yfinance)
|   |-- sentiment/      # Placeholder (Implementierung in NB02)
|   `-- utils/          # Konfigurationshelfer
|-- notebooks/
|   |-- 01_data_quality_eda.ipynb     # Datenqualität & EDA
|   |-- 02_sentiment_analysis.ipynb   # Sentiment-Inferenz
|   `-- 03_correlation_modeling.ipynb # Korrelationsanalyse
|-- data/
|   |-- raw/prices/     # prices.csv + prices_meta.json
|   `-- processed/      # Parquet-Dateien (Sentiment, Korrelationen)
|-- configs/
|   |-- tickers.yml     # Fokus-Ticker & Start-Datum
|   `-- .env            # Secrets (gitignored)
`-- report/             # LaTeX-Quellen & Abbildungen
\end{lstlisting}

\subsection{Datenfluss}

Der Informationsfluss durch das System ist strikt unidirektional und in drei
Stufen gegliedert (Abbildung~\ref{fig:datenfluss}):

\begin{figure}[H]
\centering
\begin{tabular}{ccccccccc}
\texttt{X.com} & $\longrightarrow$ & \texttt{SQLite} & $\longrightarrow$ &
\texttt{Parquet} & $\longrightarrow$ & \texttt{Ergebnis-Parquets} \\
& & $\uparrow$ & & & & $\downarrow$ \\
& & \texttt{yfinance} & & & & \texttt{Figuren / PDF}\\
\end{tabular}
\caption{Datenfluss im WebMiningProject}
\label{fig:datenfluss}
\end{figure}

\subsection{Datenhaltung}

Für persistente Speicherung werden zwei Formate eingesetzt:

\begin{itemize}[noitemsep]
  \item \textbf{SQLite} (\texttt{data/webmining.db}): Rohtweets,
        Engagement-Snapshots und Intraday-Preis-Snapshots.
        Betrieb im WAL-Modus mit aktivierten Foreign~Keys.
  \item \textbf{Parquet / CSV}: Verarbeitete Daten für Notebooks.
        Parquet bietet typsichere Spalten, schnelle Lesepfade und effiziente
        Komprimierung gegenüber CSV bei gleichzeitiger Kompatibilität mit
        pandas und PyArrow.
\end{itemize}

% =============================================================================
\section{Datenpipeline: X/Twitter-Crawler}
\label{sec:crawler}
% =============================================================================

\subsection{Technologieauswahl: Playwright statt API}

X\,(Twitter) hat im Jahr 2023 den kostenlosen API-Zugang für Drittaggregation
weitgehend eingestellt. Als Alternative wurde ein browserbasierter Crawler auf
Basis von \textbf{Playwright} (headless Chromium) implementiert. Playwright
erlaubt die Injektion echter Session-Cookies und damit die Nutzung des regulären
X-Frontends ohne kostenpflichtige API-Lizenz.

Der Crawler ist explizit für den Betrieb auf einem \textbf{Raspberry~Pi~4}
(ARM64) ausgelegt. Da Playwright keine vorgebauten ARM-Chromium-Binaries
bereitstellt, verweist die Konfiguration auf den Systemchromium
(\texttt{chromium-browser}), der via \texttt{apt} installiert wird.
Das Skript \texttt{setup\_pi.sh} automatisiert die vollständige Einrichtung
inklusive Cron-Jobs (Discovery Mo--Fr um 09:00, 12:00, 15:00~Uhr;
Polling alle 30~Minuten).

\subsection{Datenbankschema}

Die SQLite-Datenbank enthält drei Tabellen:

\begin{table}[H]
\centering
\caption{SQLite-Schema (\texttt{data/webmining.db})}
\begin{tabular}{llp{7cm}}
\toprule
\textbf{Tabelle} & \textbf{PK} & \textbf{Wesentliche Spalten} \\
\midrule
\texttt{tweets} & \texttt{tweet\_id} (TEXT) &
  \texttt{author}, \texttt{author\_followers}, \texttt{content},
  \texttt{posted\_at} (ISO~8601), \texttt{symbol},
  \texttt{likes}, \texttt{retweets} \\
\addlinespace
\texttt{tweet\_snapshots} & \texttt{id} (INT) &
  \texttt{tweet\_id} (FK), \texttt{crawled\_at} (UTC),
  \texttt{likes}, \texttt{retweets} \\
\addlinespace
\texttt{price\_snapshots} & \texttt{id} (INT) &
  \texttt{tweet\_id} (FK), \texttt{symbol},
  \texttt{timestamp} (1\,min OHLCV),
  \texttt{price}, \texttt{volume} \\
\bottomrule
\end{tabular}
\end{table}

\noindent
\texttt{store\_tweet()} nutzt den \texttt{UNIQUE}-Constraint auf
\texttt{tweet\_id} zur idempotenten Ausführung: ein \texttt{IntegrityError}
wird abgefangen und als \texttt{False} zurückgegeben (Duplikat), während eine
erfolgreiche Insertion \texttt{True} zurückgibt. Dies ermöglicht
restart-sichere Crawler-Läufe.

\subsection{Such-Query-Design}

Die Konfiguration (\texttt{x\_config.py}) definiert über 100~Suchanfragen in
sechs Sprachen (EN, DE, FR, IT, JA, SV). Jede Query kombiniert mehrere
Signaltypen:

\begin{lstlisting}[language=Python, caption={Beispiel-Queries (x\_config.py)}]
QUERIES = [
    "$TSLA OR #Tesla OR #TeslaStock",         # Cashtag + Hashtag EN
    "#Tesla #Aktie OR #TeslaAktie",           # Hashtag DE
    "Tesla OR $TSLA lang:ja",                # Japanisch
    "$BYD OR #BYDCar OR #\u6bd4\u4e9a\u8fea",          # BYD + CJK-Zeichen
    "$NIBE OR #NIBEHeatPump OR #Waermepumpe", # NIBE + Sektorterm
]
\end{lstlisting}

Server-seitig werden \texttt{min\_faves:} und \texttt{since:} in die URL
injiziert. Client-seitig filtert der Crawler nach \texttt{MIN\_LIKES\,=\,100}
und \texttt{MIN\_FOLLOWERS\,=\,100}.

\subsection{Zweistufiges Ticker-Matching}

Da Tweets keinen einheitlichen Ticker-Tag tragen, erfolgt die Zuordnung zu
Fokus-Unternehmen in zwei Passes:

\begin{enumerate}[noitemsep]
  \item \textbf{Pass~1 (direkt)}: Cashtag-Aliases werden bekannten Symbolen
        zugeordnet (z.\,B.\ \texttt{\$BYD} $\rightarrow$ \texttt{BYDDY},
        \texttt{\$VOW} $\rightarrow$ \texttt{VWAGY}).
  \item \textbf{Pass~2 (Regex)}: Tweets ohne Symbolzuordnung werden anhand
        regulärer Ausdrücke auf Firmennamen und Marken geprüft.
        Wortgrenzen wurden auf \texttt{(?<!\textbackslash w)} angepasst,
        um zusammengesetzte Hashtags (\texttt{\#TeslaStock},
        \texttt{\#NIBEHeatPump}) zu erfassen.
\end{enumerate}

\subsection{Anti-Bot-Maßnahmen und Polling-Logik}

Zur Vermeidung von Sperrungen implementiert der Crawler mehrere Strategien:

\begin{itemize}[noitemsep]
  \item \textbf{Zufällige Delays}: 3--6\,s zwischen Requests
        (\texttt{REQUEST\_DELAY\_MIN/MAX}).
  \item \textbf{Lazy Follower-Abfrage}: Follower-Zahlen werden nur für Tweets
        erhoben, die den Like-Filter bestanden haben (Profilbesuche erhöhen
        das Entdeckungsrisiko).
  \item \textbf{Adaptiver Velocity-Filter}: Nach mindestens 4~Snapshots
        (\texttt{POLL\_MIN\_SNAPSHOTS}) wird ein Tweet nur weiter verfolgt,
        wenn seine Engagement-Velocity $\geq 0{,}3$~Likes/Minute beträgt.
        Dieser Filter reduziert nutzlose Polling-Zyklen für stagnante Tweets
        erheblich und wird über eine korrelierte SQL-Subquery effizient berechnet.
\end{itemize}

Das Polling-Fenster beträgt 8~Stunden (alle 30~Minuten), was bis zu 16~Snapshots
je Tweet und damit eine verwertbare Engagement-Zeitreihe erzeugt.

\subsection{Intraday-Preis-Snapshots}

Für jeden neu gespeicherten Tweet lädt der Crawler 1-Minuten-OHLCV-Daten
im Fenster $[t_{\text{Tweet}} - 30\,\text{min},\; t_{\text{Tweet}} + 120\,\text{min}]$
herunter und speichert sie in \texttt{price\_snapshots}. Geschlossene Märkte
(Wochenenden, außerhalb NYSE-Öffnungszeit 09:30--16:00~ET) werden vorab
erkannt und übersprungen.

% =============================================================================
\section{Datenpipeline: Historische Kurspreise}
\label{sec:preise}
% =============================================================================

Der Kursdaten-Fetcher (\texttt{src/finance/price\_fetcher.py}) lädt historische
OHLCV-Tagesdaten für alle Fokus-Ticker über \textbf{yfinance}. Das Design folgt
dem Prinzip einer wachsenden, inkrementell aktualisierten CSV-Datei:

\begin{lstlisting}[language=Python, caption={Inkrementelles Laden (price\_fetcher.py)}]
# Pro Ticker nur fehlende Tage nachladen
for ticker in tickers:
    last = existing[existing["Ticker"] == ticker]["Date"].max()
    start = last + timedelta(days=1)
    new_data = yf.download(ticker, start=start, end=end,
                           auto_adjust=True, progress=False)
\end{lstlisting}

Neben der CSV wird eine JSON-Metadatendatei (\texttt{prices\_meta.json}) mit
Zeitstempel, Ticker-Liste, Zeilenanzahl, Datumsbereich und yfinance-Version
gespeichert. Ein 1{,}5-Sekunden-Delay zwischen Ticker-Downloads verhindert
Rate-Limiting bei Yahoo Finance.

Für die Analyse wird ausschließlich \texttt{Adj~Close} (split- und
dividendenbereinigt) als Kursreferenz verwendet. Kursausreißer wurden in NB01
über die IQR\,$\times$\,3-Regel je Ticker identifiziert; fehlende Handelstage
wurden mit der Bibliothek \texttt{exchange\_calendars} (NYSE-Kalender) überprüft.

% =============================================================================
\section{Sentiment-Analyse-Pipeline}
\label{sec:sentiment}
% =============================================================================

\subsection{Modell und Vorverarbeitung}

Die Sentiment-Inferenz erfolgt in NB02 mit dem vortrainierten Modell
\texttt{cardiffnlp/twitter-xlm-roberta-base-sentiment}
\autocite{loureiro2022timelms} über die Hugging~Face
Transformers-Bibliothek. Die Wahl dieses Modells begründet sich in seiner
Mehrsprachigkeit: Der Datensatz enthält Tweets auf Englisch, Deutsch,
Französisch und Japanisch, die ohne Übersetzungsumweg verarbeitet werden können.

Vor der Inferenz werden Texte bereinigt:

\begin{lstlisting}[language=Python, caption={Textbereinigung vor Modell-Inferenz}]
import re

def clean_text(text: str) -> str:
    text = re.sub(r"https?://\S+", "", text)   # URLs entfernen
    text = re.sub(r"[$@]\w+", "", text)        # Cashtags & Mentions
    return text.strip()

# Nur Tweets mit mind. 20 verbleibenden Zeichen inferieren
df["has_content"] = df["content"].map(clean_text).str.len() >= 20
\end{lstlisting}

Das Modell liefert Wahrscheinlichkeiten für drei Klassen; der kontinuierliche
Compound-Score ist:
\begin{equation}
\texttt{xlm\_compound} = P(\text{positive}) - P(\text{negative}) \in [-1,\, 1]
\end{equation}

Die Inferenz läuft im Batch-Modus (\texttt{BATCH\_SIZE=32}) mit optionaler
GPU-Beschleunigung. Ein Caching-Mechanismus (\texttt{FORCE\_RERUN=False})
verhindert redundante GPU-Nutzung bei unveränderten Eingabedaten.

\subsection{Tägliche Aggregation und Follower-Gewichtung}

Tweet-Level-Scores werden je Fokus-Ticker und Kalendertag aggregiert.
Das primäre Analysesignal ist der follower-gewichtete Compound-Score:

\begin{equation}
\texttt{weighted\_xlm} = \frac{\sum_{i} f_i \cdot c_i}{\sum_{i} f_i}
\end{equation}

wobei $f_i$ die Follower-Zahl und $c_i$ der Compound-Score von Tweet~$i$.
Tage, an denen alle Autoren null Follower haben, erhalten \texttt{NaN}
(werden beim Inner-Join in NB03 implizit ausgeschlossen).
Die Ausgabedatei \texttt{daily\_sentiment\_original.parquet}
enthält 508~Zeilen über 12~Ticker.

\subsection{Retweet-Separation}

Retweets (erkennbar am Präfix \texttt{"RT @"}) werden in separate
Parquet-Dateien ausgelagert (\texttt{*\_retweets.parquet}), um das
Primärsignal nicht durch bloße Multiplikation originärer Meinungen zu
verzerren. Im vorliegenden Datensatz enthielten alle gecrawlten Inhalte
originären Text; die Separation ist methodisch korrekt implementiert und
für Folgestudien wiederverwendbar.

% =============================================================================
\section{Korrelationsmodellierung}
\label{sec:korrelation}
% =============================================================================

\subsection{Daten-Merge und Overlap-Report}

NB03 führt Sentiment-Daten und Kurspreise über einen Inner-Join auf
\texttt{(focus\_ticker, date)} zusammen. Statt Ticker mit zu wenig Daten
stillschweigend auszuschließen, erzeugt das Notebook einen
\textbf{Overlap-Report} mit methoden-spezifischen Eligibility-Flags, der als
\texttt{overlap\_report.parquet} persistent gespeichert wird:

\begin{table}[H]
\centering
\caption{Eligibility-Schwellwerte im Overlap-Report}
\begin{tabular}{ll}
\toprule
\textbf{Methode} & \textbf{Mindestvoraussetzung} \\
\midrule
Pearson/Spearman     & $\geq 5$ gepaarte Beobachtungen je Lag \\
Rolling-Korrelation  & $\geq 30$ überlappende Tage \\
Granger-Test         & $\geq 30$ Beobachtungen \\
Event-Study          & $\geq 20$ Basistage \emph{und} $\geq 5$ Ereignisse je Richtung \\
\bottomrule
\end{tabular}
\end{table}

\subsection{Methodische Implementierungsdetails}

\paragraph{Multiple Testing (BH-FDR).}
Da Pearson/Spearman-Tests für bis zu 10~Ticker $\times$ 4~Lags = 40 simultane
Hypothesen durchgeführt werden, korrigiert das Notebook alle $p$-Werte
mit der Benjamini-Hochberg-FDR-Methode:

\begin{lstlisting}[language=Python, caption={BH-Korrektur via statsmodels}]
from statsmodels.stats.multitest import multipletests

_, p_adj, _, _ = multipletests(corr_df["pearson_p"].values, method="fdr_bh")
corr_df["pearson_p_adj"] = p_adj
corr_df["pearson_significant"] = p_adj < 0.05
\end{lstlisting}

\paragraph{Konfidenzintervalle (Fisher-z).}
Zu jedem Pearson-$r$ wird ein 95\,\%-KI über die Fisher-$z$-Transformation
berechnet ($\mathrm{SE}_z = (n-3)^{-1/2}$), das die Schätzunsicherheit bei
kleinen Stichproben transparent macht.

\paragraph{Stationaritätsprüfung (ADF).}
Vor den Granger-Tests werden alle Eingangszeitreihen mit \texttt{adfuller()}
aus \texttt{statsmodels} auf Einheitswurzeln geprüft.
Alle vier getesteten Ticker waren auf beiden Reihen stationär
($p < 0{,}001$) -- die Voraussetzung war erfüllt.

\paragraph{Event-Study-Benchmark.}
Als Markt-Benchmark für abnormale Renditen wird der S\&P~500 (\texttt{\^{}GSPC})
via yfinance heruntergeladen, anstatt den internen Focus-Ticker-Mittelwert zu
verwenden, der zu einem zirkulären Vergleich geführt hätte:

\begin{lstlisting}[language=Python, caption={Externer S\&P-500-Benchmark}]
spx_raw = yf.download("^GSPC", start=spx_start, end=spx_end,
                       auto_adjust=True, progress=False)
market_ret = spx_raw["Close"].pct_change().rename("market_return")
merged2["abnormal_return"] = merged2["return"] - merged2["market_return"]
\end{lstlisting}

\subsection{Ausgabeformate}

\begin{table}[H]
\centering
\caption{Ergebnis-Parquets aus NB03 (\texttt{data/processed/})}
\begin{tabular}{ll}
\toprule
\textbf{Datei} & \textbf{Inhalt} \\
\midrule
\texttt{correlation\_results.parquet}  &
  $r$, KI, $p_{\mathrm{roh}}$, $p_{\mathrm{adj}}$, $n$ je (Ticker, Lag) \\
\texttt{granger\_results.parquet}      &
  $p$-Werte (F-Test) je (Ticker, Richtung, Lag) \\
\texttt{event\_study\_results.parquet} &
  Mittlere abnormale Rendite je (Ticker, Richtung, $t$) \\
\texttt{overlap\_report.parquet}       &
  Coverage und Eligibility-Flags je Fokus-Ticker \\
\bottomrule
\end{tabular}
\end{table}

% =============================================================================
\section{Konfiguration, Tests und Reproduzierbarkeit}
\label{sec:konfiguration}
% =============================================================================

\subsection{Konfigurationsmanagement}

Secrets (X-Session-Cookies) werden ausschließlich über
\texttt{configs/.env} (gitignored) bereitgestellt; eine kommentierte
Vorlagedatei \texttt{configs/.env.example} ist versioniert und dokumentiert
alle erwarteten Variablen. Analyse-Parameter sind zentral definiert:

\begin{itemize}[noitemsep]
  \item \texttt{configs/tickers.yml}: 13~Fokus-Ticker + \texttt{start\_date}
  \item \texttt{src/scraping/x\_config.py}: alle Crawler-Parameter als
        benannte Konstanten (\texttt{MIN\_LIKES}, \texttt{POLL\_INTERVAL\_MINUTES}, \ldots)
\end{itemize}

Jedes Notebook definiert Schwellwerte am Anfang explizit, sodass
Reproduzierbarkeit ohne Kenntnis externer Konfigurationsdateien gegeben ist.

\subsection{Automatisierte Tests}

\begin{table}[H]
\centering
\caption{pytest-Testdateien in \texttt{tests/}}
\begin{tabular}{ll}
\toprule
\textbf{Datei} & \textbf{Geprüfte Komponente} \\
\midrule
\texttt{test\_db.py}             &
  SQLite-Schema, \texttt{store\_tweet()}, Duplikaterkennung \\
\texttt{test\_price\_fetcher.py} &
  Inkrementelles Laden, Merge-Logik, Metadaten-JSON \\
\texttt{test\_x\_crawler.py}     &
  Parser-Funktionen (\texttt{\_parse\_stat()}, Filterlogik) \\
\bottomrule
\end{tabular}
\end{table}

\subsection{Notebook-Konventionen}

\begin{itemize}[noitemsep]
  \item Outputs werden vor Commits geleert (\emph{Clear All Outputs}).
  \item Zufallsprozesse nutzen feste Seeds (\texttt{seed=42}).
  \item Figuren werden via \texttt{\_save\_fig()} automatisch nach
        \texttt{report/figures/} exportiert (\texttt{FIG\_DPI=300}).
  \item Parquet-Artefakte am Ende jedes Notebooks gespeichert;
        Folgenotebooks lesen ausschließlich diese Dateien
        (keine direkten SQLite-Abhängigkeiten in NB02/NB03).
\end{itemize}

% =============================================================================
\section{Fazit}
\label{sec:fazit}
% =============================================================================

Die implementierte Pipeline deckt den vollständigen Datenlebenszyklus von der
Erhebung bis zur statistischen Auswertung ab. Besonderes Augenmerk wurde auf
drei Aspekte gelegt:

\begin{enumerate}[noitemsep]
  \item \textbf{Modularität}: Klare Trennung zwischen Scraping, Preisdaten und
        Analyse; jede Komponente ist unabhängig testbar und austauschbar.
  \item \textbf{Reproduzierbarkeit}: Parquet-Artefakte als unveränderliche
        Zwischenstände, explizite Parameter, versionierte Konfigurationsvorlagen.
  \item \textbf{Methodische Korrektheit}: BH-FDR-Korrektur für Multiple~Testing,
        ADF-Stationaritätstests vor Granger-Kausalitätstests,
        externer S\&P~500-Benchmark für die Event-Study.
\end{enumerate}

Erweiterungspotenzial besteht insbesondere in drei Bereichen:
(1)~Intraday-Analyse auf Basis der bereits gespeicherten
\texttt{price\_snapshots}-Daten,
(2)~Log-Returns für normalverteilungsnahere Residuenverteilungen im
Granger-Modell sowie
(3)~Kreuzvalidierung zur Reduktion des Data-Snooping-Bias bei der Lag-Selektion.

\printbibliography[heading=bibintoc]

\end{document}
